\documentclass[11pt]{article}
\usepackage[T1]{fontenc}
\usepackage{lmodern,microtype,amsmath,amssymb,amsthm,booktabs,longtable,array}
\usepackage[margin=1in]{geometry}
\usepackage{xcolor,xurl}
% Shared current-manuscript references for both companion documents.
% Rebuild hi_cft_c8_proof.tex before either companion after manuscript edits.
\RequirePackage{xr-hyper}
\RequirePackage[colorlinks=true,linkcolor=blue!50!black,urlcolor=blue!50!black,
  filecolor=blue!50!black]{hyperref}
\IfFileExists{hi_cft_c8_proof.aux}{}{%
  \PackageError{manuscript-refs}{Missing hi_cft_c8_proof.aux}%
    {Compile hi_cft_c8_proof.tex before compiling the companion documents.}}
\externaldocument[ms-]{hi_cft_c8_proof}[hi_cft_c8_proof.pdf]
\newcommand{\msref}[1]{\ref{ms-#1}}
\newcommand{\mspageref}[1]{\pageref{ms-#1}}
\newcommand{\msnameref}[1]{\nameref{ms-#1}}

\newcommand{\QQ}{\mathbb Q}\newcommand{\CC}{\mathbb C}\newcommand{\ZZ}{\mathbb Z}
\newcommand{\Span}{\operatorname{span}}\newcommand{\Rep}{\operatorname{Rep}}
\newcommand{\HH}{\mathcal H}\newcommand{\one}{\mathbf1}
\newcommand{\artifact}[1]{\path{#1}}
\newcolumntype{P}[1]{>{\raggedright\arraybackslash}p{#1}}
\newtheorem{proposition}{Proposition}[section]
\newtheorem{lemma}[proposition]{Lemma}
\newtheorem{theorem}[proposition]{Theorem}
\newtheorem{corollary}[proposition]{Corollary}
\theoremstyle{remark}\newtheorem{remark}[proposition]{Remark}
\setlength{\emergencystretch}{2em}
\title{Human verification of the finite inputs\\to the $c=8$ Haagerup construction}
\author{Companion to the revised research manuscript}
\date{27 September 2026}
\begin{document}
\maketitle
\begin{abstract}
The construction uses exact calculations, but a list of successful jobs
does not explain what a mathematician should inspect. This guide replaces
that interface with explicit scalar equations, sparse field coordinates,
membership witnesses and a small independent checker. It separates three
tasks: checking a displayed arithmetic identity, checking a finite vector
identity in the lattice model, and proving the theorem that turns those
identities into an all-weight or categorical conclusion. The first two
are made more accessible here; the third remains in the manuscript.
\end{abstract}
\tableofcontents

\section{How to audit the inputs}
The reference manuscript is the expanded physics revision of
\href{hi_cft_c8_proof.pdf}{\artifact{hi_cft_c8_proof.pdf}}.
References explicitly marked ``manuscript'' use its current section,
appendix and proposition numbers. Their linked numbers and pages are
imported from the manuscript's own labels; unqualified references are
internal to this guide. Keep the two PDFs in the same directory to follow
the links. The original certificates remain unchanged.
This companion makes selected inputs inspectable without
asking the reader to trust a hash, a rank counter, or a printed
\texttt{true} value as a mathematical proof.

A useful audit proceeds in three passes.
\begin{enumerate}
\item Check the small equations below: the current levels, the three
momenta, the stress projection matrix, the explicit parafermion rows,
and the short Gate~5 relation. These are literal formulas, with all
coefficient conventions fixed.
\item For larger finite equalities, inspect a witness rather than repeat
the search that discovered it. The necessary format is an equation
$v=\sum_j c_jw_j$, the coordinates of each vector, and a recipe identifying
$w_j$ with the claimed lattice product. A rank lower bound needs a
nonzero minor or an elimination transcript; a rank upper bound needs
spanning expressions. These are different obligations.
\item Follow the printed implications in manuscript
Appendices~\msref{sec:finite}--\msref{sec:analytic-full-theory}.
An exact residual of zero proves its finite identity. It does not by
itself prove strong generation, regularity, module exhaustion, or
complete unitarity.
\end{enumerate}

\begin{center}\small
\begin{tabular}{@{}P{.25\linewidth}P{.65\linewidth}@{}}
\toprule Input & Human-facing evidence supplied here\\\midrule
Lattice, currents, $u$ & An eight-dimensional integer calculation and
explicit coefficient extraction; all $72$ current singular products
are independently checked.\\
Positive stress identity & Four orthogonal stress norms and eight actual
overlaps; the resulting vector identity is independently checked.\\
Low-grade fields & All $337$ monomials of eight relevant grade-two fields
are exported in plain tab-separated form.\\
Short generation & A $120$-row construction graph, exact elimination
identities and expansions of all ten target fields in two bases.\\
Factor membership & A $33$-word Potts expression and a $42$-word
parafermion expression, with full coordinates and an independent check.\\
Gates~4 and~5 & Denominator-cleared equations, named normal words,
nonzero scalar witnesses and a precise residual format.\\
Embedded factors and tops & Full first-row OPE formulas, charge and
spectral tables, an $80$-row cyclic basis and the precise distinction
between finite matrix evidence and its mode interpretation.\\
\bottomrule
\end{tabular}\end{center}

\subsection{Navigation in the expanded manuscript}
The following references locate the current arguments supported by
these finite inputs. Page numbers refer to the manuscript, not to this
guide. The Lean companion supplies a section-by-section assessment of
the remaining formalization work.
\begin{longtable}{@{}P{.28\linewidth}P{.60\linewidth}@{}}
\toprule Audit topic & Current manuscript location\\\midrule
\endhead
Lattice and initial spaces & Appendix~\msref{sec:construction},
p.~\mspageref{sec:construction}; Lemma~\msref{lem:levels} and
Certificate~\msref{cert:initial}.\\
All-weight closure & Appendix~\msref{sec:finite},
p.~\mspageref{sec:finite}; Theorem~\msref{thm:recursive} and
Certificate~\msref{cert:lower}.\\
Gates and properness & Appendices~\msref{sec:gate4}--\msref{sec:properness},
pp.~\mspageref{sec:gate4}--\mspageref{sec:properness};
Lemmas~\msref{lem:singlet}, \msref{lem:bootstrap}, and \msref{lem:jacobi}.\\
Embedded factors and stresses & Appendix~\msref{sec:regularity},
p.~\mspageref{sec:regularity}; Certificate~\msref{cert:structural-inputs},
Proposition~\msref{prop:pf13-finite-reconstruction}, and
Theorem~\msref{thm:positive-stress-rationality}.\\
Branching and characters & Appendices~\msref{sec:branching}--\msref{sec:characters},
pp.~\mspageref{sec:branching}--\mspageref{sec:characters}.\\
Analytic and categorical conclusions & Appendices~\msref{sec:analytic-unitarity}--\msref{sec:analytic-full-theory},
pp.~\mspageref{sec:analytic-unitarity}--\mspageref{sec:analytic-full-theory}.\\
Physical observables & Sections~\msref{sec:physics-spectrum}--\msref{sec:physics-data-guide},
pp.~\mspageref{sec:physics-spectrum}--\mspageref{sec:physics-data-guide};
the data guide is Section~\msref{sec:physics-data-guide}.\\
Evidence and formalization scope & Appendices~\msref{sec:ledger}--\msref{sec:compilation},
pp.~\mspageref{sec:ledger}--\mspageref{sec:compilation}.\\
\bottomrule
\end{longtable}

\section{The small lattice calculation}
Put $s^2=13$, $t^2=-3$, and use the coefficient field
$K=\QQ(s,t)$ with ordered basis $(1,s,t,st)$. Complex conjugation fixes
$s$ and negates $t$. If $x=(a,b,c,d)$ and $y=(e,f,g,h)$, then
\[
xy=(ae+13bf-3cg-39dh,\quad af+be-3ch-3dg,
\quad ag+ce+13bh+13df,\quad ah+de+bg+cf).
\]
Thus every vector equality is exactly four rational equalities per
lattice monomial. No numerical recognition tolerance occurs in the
checker. The historical field file also contains floating-point discovery
data; only its exact rational coefficient arrays are read.

The lattice Gram matrix is
\[
G=\begin{pmatrix}
2&-1&0&0&0&0&0&0\\-1&2&-1&0&0&0&0&0\\
0&-1&2&-1&0&-1&0&0\\0&0&-1&2&-1&0&0&0\\
0&0&0&-1&2&0&0&0\\0&0&-1&0&0&2&0&0\\
0&0&0&0&0&0&2&-1\\0&0&0&0&0&0&-1&2
\end{pmatrix}.
\]
Let $\ell=(-1,-2,-3,-2,-1,-2,0,0)$ and
$r=(5,8,9,8,5,0,2,2)$. Two matrix multiplications give
\[
G\ell=(0,0,0,0,0,-1,0,0),\qquad
Gr=(2,2,2,2,2,-9,2,2).
\]
Consequently $\ell^TG\ell=2$, $r^TGr=78$, and
$\ell^TG\alpha_i=0$ for every $i$ with $r_i\ne0$.
These are the geometric ingredients of levels $1$ and $39$.
Explicitly, with $I=\{0,1,2,3,4,6,7\}$, the six currents are
\[
\begin{aligned}
H_L&=\ell(-1)\one,& E_L&=e^\ell,& F_L&=-e^{-\ell},\\
H_R&=r(-1)\one,& E_R&=\sum_{i\in I}e^{\alpha_i},&
F_R&=-\sum_{i\in I}r_i e^{-\alpha_i}.
\end{aligned}
\]

For signs use
\[
\epsilon(\alpha,\beta)=(-1)^{\sum_i\alpha_i\beta_i+
 \sum_{i>j}\alpha_i\beta_jG_{ij}}.
\]
For example $\epsilon(\alpha_i,-\alpha_i)=-1$, so
$e^{\alpha_i}_0e^{-\alpha_i}=-\alpha_i(-1)\one$ and
$e^{\alpha_i}_1e^{-\alpha_i}=-\one$.
The minus signs in the lowering currents therefore give
\[
(E_R)_0F_R=\sum_i r_i\alpha_i(-1)\one=H_R,
\qquad (E_R)_1F_R=\sum_i r_i\one=39\one.
\]
Together with the Heisenberg brackets this checks the complete affine
relations. The reader can derive the mixed vanishing from the displayed
orthogonality. The independent checker evaluates all ordered pairs of
the six currents at modes $0$ and $1$, rather than checking only these
two examples. Higher singular modes vanish by weight.

\subsection{The three-term vector without a mode engine}
\label{hr:three-term-vector}
Write $u=e^a+(s+3t)e^b/20+(s-3t)e^c/20$, with
\[
\begin{split}
a&=(0,0,0,0,0,-1,1,1),\\
b&=(0,0,0,1,1,-1,0,0),\\
c&=(1,1,0,0,0,-1,0,0).
\end{split}
\]
The following table can be checked directly using $G$.
\begin{center}\small
\begin{tabular}{@{}ccccc@{}}\toprule
$\alpha$ & $G\alpha$ & $\alpha^TG\alpha$ & $\ell^TG\alpha$ & $r^TG\alpha$\\\midrule
$a$ & $(0,0,1,0,0,-2,1,1)$ & $4$ & $1$ & $13$\\
$b$ & $(0,0,0,1,1,-2,0,0)$ & $4$ & $1$ & $13$\\
$c$ & $(1,1,0,0,0,-2,0,0)$ & $4$ & $1$ & $13$\\\bottomrule
\end{tabular}\end{center}
For each summand the inner product with a raising root is nonnegative.
The coefficient of $z^{-1}$ in
\[
Y(e^\alpha,z)e^\beta
=\epsilon(\alpha,\beta)z^{\alpha^TG\beta}
\exp\left(\sum_{n>0}\frac{\alpha(-n)}n z^n\right)e^{\alpha+\beta}
\]
therefore vanishes for every raising current. The absolute inner
product with a root occurring in either current triple is at most one,
so positive current modes also vanish. The Cartan positive modes kill
oscillator-free exponentials. Thus $u$ is an affine primary highest
vector of current dimensions $(2,14)$; this check uses no cancellation
between its three coefficients.

Set $\widetilde r=(5,8,9,8,5,1,2,2)$ and
$q(\alpha)=\prod_i\widetilde r_i^{\alpha_i}$. In the transported
positive form, $\|e^\alpha\|^2=q(\alpha)$. Here
$q(a)=4$ and $q(b)=q(c)=40$, while
\[
\left|\frac{s\pm3t}{20}\right|^2=\frac{13+27}{400}=\frac1{10}.
\]
The three orthogonal summands each have norm squared four, and
$\|u\|^2=12$. This explains the equal-magnitude phase convention
without confusing the transported PCT with ordinary conjugation.

\section{A four-column check of the positive stress identity}
\label{hr:positive-stress-check}
The ordered axes are $w=(\omega_L,\omega_R,e,f)$. Their Gram matrix
is diagonal with diagonal
\[
N=\left(\frac12,\frac{117}{82},\frac{1200}{1271},\frac{35}{31}\right).
\]
The entries can be checked in two independent ways: from the stated
commuting Virasoro structures, or from the sparse fields and the lattice
Fock inner product. The new checker uses the latter, so it does not
assume the claimed central charges to compute the pairings.

For precision, if $\alpha=\beta$, the pairing of two oscillator
monomials with $r$ factors is
\[
\left\langle\prod_{k=1}^r\alpha_{i_k}(-n_k)e^\alpha,
\prod_{k=1}^r\alpha_{j_k}(-m_k)e^\alpha\right\rangle_g
=q(\alpha)\sum_{\sigma\in S_r}
\prod_{k=1}^r\delta_{n_k,m_{\sigma(k)}}\,n_kG_{i_kj_{\sigma(k)}}.
\]
Different momenta or different numbers of oscillator factors are
orthogonal. Coefficients in the first argument are conjugated. In
weight two there are at most two oscillator factors, so the permutation
sum has at most two terms. For example the pairing of
$\alpha_i(-1)\alpha_j(-1)\one$ with
$\alpha_k(-1)\alpha_l(-1)\one$ is
$G_{ik}G_{jl}+G_{il}G_{jk}$, while the pairing of the two mode-two
oscillators is $2G_{ij}$.

The freshly recomputed overlaps and normalized projections are
\begin{center}
\begin{tabular}{@{}ccccc@{}}\toprule
 & $\omega_L$ & $\omega_R$ & $e$ & $f$\\\midrule
$\langle w_i,p\rangle_g$ & $0$ & $13/205$ & $108/1271$ & $39/155$\\
$\langle w_i,z\rangle_g$ & $0$ & $21/205$ & $780/1271$ & $13/155$\\\midrule
$\langle w_i,p\rangle_g/N_i$ & $0$ & $2/45$ & $9/100$ & $39/175$\\
$\langle w_i,z\rangle_g/N_i$ & $0$ & $14/195$ & $13/20$ & $13/175$\\\bottomrule
\end{tabular}\end{center}
The squared norms of the actual saved fields are
$\langle p,p\rangle_g=2/5$ and $\langle z,z\rangle_g=4/5$.
These norm checks alone do not verify their Virasoro OPEs.

The four columns reduce the positive identity to
\[
\begin{split}
1&=1,\\
1&=\frac{3403}{4563}+\frac{25}{6}\frac{2}{45}
                         +\frac{25}{26}\frac{14}{195}
   =\frac{3403+845+315}{4563},\\
1&=\frac{25}{6}\frac9{100}+\frac{25}{26}\frac{13}{20}
   =\frac38+\frac58,\\
1&=\frac{25}{6}\frac{39}{175}+\frac{25}{26}\frac{13}{175}
   =\frac{13}{14}+\frac1{14}.
\end{split}
\]
Thus the exact vector identity
\[
\omega_B=\omega_L+\frac{3403}{4563}\omega_R
                      +\frac{25}{6}\overline p+\frac{25}{26}\overline z
\]
has a small inspectable arithmetic certificate. The identification of
the orthogonal projections with Haar averages still uses the separately
certified statement that these four axes exhaust the current invariants
in the actual $S_2$. The checker verifies the projected vector equality;
it does not replace that representation-theoretic input.

All $337$ monomials of $u,\omega_B,\omega_L,\omega_R,e,f,p,z$ are
printed as machine-independent rows in
\artifact{c8_reviewability_20260927/grade2_fields.tsv}. Each row has
the field name, eight momentum coordinates, the oscillator indices,
and the four rational coefficients. An oscillator index $8(n-1)+i$
means $\alpha_i(-n)$. Repeated indices are repeated commuting creation
operators, with no divided-power convention.

\section{Finite closure: identities that a reader can inspect}
\label{sec:reader-finite}

The finite closure argument in Appendices~\msref{sec:construction}--\msref{sec:properness} of the
\href{hi_cft_c8_proof.pdf}{manuscript} has three different kinds of
evidence.  A sparse identity gives an equality of actual vectors.  A
nonzero minor proves a positive rank assertion.  A file hash identifies
the source of either witness.  Only the first two supply mathematical
information; the third prevents a change of input from going unnoticed.
This section makes the distinction concrete and prints the small
identities responsible for the final Jacobi step.

\subsection{How to read a coefficient and a row}

Put $F=\QQ(s,t)$, where $s^2=13$ and $t^2=-3$.  The four rational
strings $[a,b,c,d]$ in a record mean
\[
                  a+bs+ct+dst.
\]
The order is significant.  For example, the coefficient
$[0,-54087/138650,20899/138650,0]$ is
$(-54087s+20899t)/138650$.  Multiplication can be checked using
only rational arithmetic: if $x=(x_0,x_1,x_2,x_3)$ and
$y=(y_0,y_1,y_2,y_3)$, the product has coordinates
\begin{align*}
z_0&=x_0y_0+13x_1y_1-3x_2y_2-39x_3y_3,\\
z_1&=x_0y_1+x_1y_0-3(x_2y_3+x_3y_2),\\
z_2&=x_0y_2+x_2y_0+13(x_1y_3+x_3y_1),\\
z_3&=x_0y_3+x_3y_0+x_1y_2+x_2y_1.
\end{align*}
In particular, zero testing requires four rational zero tests; it
does not involve a numerical tolerance or a choice of decimal roots.

A vector is a finite list of triples
\[
       [\,\lambda,\ \mathcal O,\ [a,b,c,d]\,].
\]
Here $\lambda\in\ZZ^8$ gives the momentum in the ordered simple-root
basis, and a nonnegative integer $8(n-1)+i$ in the oscillator list
$\mathcal O$ means $\alpha_i(-n)$, with $0\leq i<8$.
The triple denotes the corresponding oscillator monomial applied to
$e^\lambda$, with coefficient $a+bs+ct+dst$.
Thus the assertion $v=\sum_jc_jw_j$ means that, for every momentum
and oscillator monomial, all four coordinates of
$v-\sum_jc_jw_j$ vanish.  Checking a pivot coordinate alone is
insufficient.

The grade-three checkpoint also records how its normalized vectors
were made.  If a row has recipe $v$, elimination list
$[(d_1,r_1),\ldots,(d_m,r_m)]$, and scale $\gamma$, its meaning is
\begin{equation}\label{eq:reader-row-recipe}
       r=\gamma\left(v-\sum_{j=1}^m d_jr_j\right).
\end{equation}
The subtraction precedes the scaling.  All referenced rows are earlier
rows.  The leading coefficient of $r$ is one, and its leading state
is different from those already inserted in the same Cartan block.
These data give both a construction by actual modes and an exact
linear-independence witness.  The recipes are not merely names for
vectors obtained by approximate diagonalization.

\subsection{The short generators really produce the saved fields}

Let $K_{\rm cur}$ denote the six-dimensional current space, and let
$u$ be the displayed three-vertex field in
Section~\ref{hr:three-term-vector}.
There is a small, deterministic calculation connecting the definition
$U=\langle K_{\rm cur},\omega_B,u\rangle$ to the much larger saved
presentation.  Start with the span of
\[
 \{x_{-2}\one,\ x_{-1}y:x,y\in K_{\rm cur}\}
                   \cup\{\omega_B,u\}.
\]
Its exact dimension is $24$.  Keep a queue of newly accepted independent
vectors.  For each vector $v$ removed from the queue, insert the five
vectors
\[
       (E_L)_0v,\quad(F_L)_0v,\quad(E_R)_0v,\quad(F_R)_0v,
                          \quad u_1v
\]
by exact sparse elimination, adding a new queue entry only when the
remainder is nonzero.  The queue empties after $120$ accepted vectors
have been processed.  Every operation preserves membership in the
already defined VOA $U$, so the resulting space lies in $U_2$.
Its dimension is $120$, and direct exact reduction puts every saved
highest vector and the retained neutral stress in this span.

Independently, the saved $S_2$ has the decomposition by current types
\[
 \dim S_2=22+2+(5+9+9+13+4+12+16+28)=120.
\]
The $22$ affine vacuum descendants, the two neutral fields, and the
eight indicated multiplets are all in the queue-generated space;
dimension comparison identifies that space with $S_2$.
This establishes $S_2\subset U_2$ before finite closure is used.
It does not yet assert $S_2=U_2$; that conclusion uses the later
strong-generation argument.

The complete short algorithm is
\artifact{hi_cft_20260923_c8_minimal_generators.py}, and its fresh
exact replay is
\artifact{c8_revision_20260927/minimal_generators_replay.json}.
The new witness
\artifact{c8_reviewability_20260927/minimal_generation_witness.json}
records all $120$ accepted raw vectors, their $120$ normalized
remainders, every elimination and scale, and the expansions of
the normalized vectors in the raw vectors.  It also records the
eight highest fields, the retained neutral stress, and $\omega_B$
in both bases.  The adjacent
\artifact{minimal_generation_witness.recipes.tsv} and
\artifact{minimal_generation_witness.target_coefficients.tsv}
give readable recipe and coordinate tables.  There were $644$
candidate insertions and $8386$ retained sparse entries in the
generation rows; the verifier counts $8442$ after including the
input-field coordinates.

The independent verification record
\artifact{c8_reviewability_20260927/minimal_generation_checked.json}
checks $120$ unit-triangular normalization identities, $120$
expansions in the raw vectors, and all $20$ target-coordinate
identities.  It also checks the six current formulas, the ambient
stress and $u$, the earlier-parent condition in the recipe graph,
and the source bindings.  Thus the coefficient identities and
the rank-$120$ witness are available explicitly.  The exporter
recomputed the mode products in $1.81$ seconds; the independent
coordinate verifier took $1.02$ seconds, with peak RSS $35900$ KiB.
The latter does not recompute a raw vector from its named mode
recipe.  That remaining interpretation is supplied by the
exporter's actual mode calculation and the lattice-mode engine.

\subsection{Gate 4: twelve projected residuals and a finite rank argument}

Write $P=S_3/N_3$.  The constructed spaces have exact dimensions
\[
           \dim S_3=956,\qquad \dim N_3=634,
                   \qquad \dim P=322.
\]
The Gate 4 conclusion is $(S_2)_0S_3\subset W_4=N_4+P_4$.
It is reduced to nine highest left inputs, each with $322$ right
inputs.  There are $2145$ original exact seed relations and $118$
additional direct relations.  The last twelve targets have weight
$(0,0)$ and are the following.  The integer in the first column is
the primary key in the projected exact database.
\begin{center}
\begin{tabular}{rcl@{\qquad}rcl}
\toprule
ID&Left&Right&ID&Left&Right\\
\midrule
0&$a_0$&\texttt{r3\_000950}&6&$a_4$&\texttt{r3\_000873}\\
1&$a_1$&\texttt{r3\_000916}&7&$a_5$&\texttt{r3\_000879}\\
2&$a_1$&\texttt{r3\_000920}&8&$a_5$&\texttt{r3\_000897}\\
3&$a_2$&\texttt{r3\_000916}&9&$a_6$&\texttt{r3\_000882}\\
4&$a_2$&\texttt{r3\_000920}&10&$a_6$&\texttt{r3\_000898}\\
5&$a_3$&\texttt{r3\_000922}&11&$a_7$&\texttt{r3\_000902}\\
\bottomrule
\end{tabular}
\end{center}
The fields $a_0,\ldots,a_7$ are the eight charged highest vectors in
their saved order.  In $(a_i)_0p$, the first subscript labels a field
and the second denotes its mode.

Here is the projector used in all twelve records, with every
denominator exposed:
\begin{align*}
 \Pi={}&\left(1-\frac{F_LE_L}{2}\right)
         \left(1-\frac{F_LE_L}{6}\right)
         \prod_{d\in D}\left(1-\frac{F_RE_R}{d}\right),\\
 D={}&\{2,6,12,20,30,42,56,72,90,110,132,156\}.
\end{align*}
The symbols in this formula are current zero modes.  On the
weight-zero line of a highest-weight-$j$ $\mathfrak{sl}_2$ module,
$FE$ acts by $j(j+2)/4$.  The current bounds at grade four are
$j_L\leq4$, $j_R\leq24$, so this polynomial kills every nontrivial
joint type in the weight-$(0,0)$ subspace and is one on the joint
singlets.

For each target $t$, a record supplies rational quadruples
$c_0,\ldots,c_{20}$, with unused coefficients omitted, and asserts
\begin{equation}\label{eq:reader-projected-residual}
           \Pi t-\sum_{j=0}^{20}c_j\Pi n_j=0,
                         \qquad n_j\in N_4.
\end{equation}
The $n_j$ are actual normal words.  Their allowed recipes are
$x_{-1}p$, $Tp$, or $a_{-1}b$, followed if necessary by current
lowering.  For example, the database recipe
\texttt{normal\_descendant}, with current index $i<6$ and input $p$,
means $(x_i)_{-1}p$; current index $6$ means $(\omega_B)_0p=Tp$.
A \texttt{current\_lowering} recipe points to an earlier normal
column and applies $(F_L)_0$ or $(F_R)_0$.

The new packet
\artifact{c8_reviewability_20260927/finite_witnesses.json}
prints the coefficient list of each of the twelve equations, the
correspondence from projected columns to original normal columns,
and the recursive normal-word recipes.  In the original exact
database, projected target $i$ is stored as \texttt{target:$i$};
projected normal column $j$ is stored as \texttt{grade5:$j$}.
The latter is a historical storage prefix: these vectors have
conformal weight \emph{four}.

Four small entries of target $11$ illustrate the complete convention.
The target is $t=(a_7)_0\texttt{r3\_000902}$, and its recorded
coefficient $c_j$ multiplies $\Pi n_j$ as follows:
\begin{center}
\begin{tabular}{rll}
\toprule
$j$&Actual normal word $n_j$&Coefficient $c_j$\\
\midrule
2&$(H_L)_{-1}\texttt{r3\_000022}$&$(169-3st)/2240$\\
3&$(H_L)_{-1}\texttt{r3\_000023}$&$(-129s+117t)/5600$\\
14&$T\texttt{r3\_000536}$&$197/6720+st/560$\\
15&$T\texttt{r3\_000537}$&$-s/1400-39t/2800$\\
\bottomrule
\end{tabular}
\end{center}
These are four entries of the $21$-term relation, not a replacement
for its other terms.  The packet retains all coefficients, including
the longer fractions; the checker uses the entire relation.

The independent reader
\artifact{c8_reviewability_20260927/finite_witnesses.py}
uses Python's standard-library rational arithmetic to recalculate
all four coordinates of every residual in
\eqref{eq:reader-projected-residual}.  It imports neither the field
engine nor a vertex-mode routine from the research scripts.  This
is a full replay of these twelve \emph{saved-vector equalities}.
The recorded run passed in $21.15$ seconds with peak process RSS
$39512$ KiB.
The construction of the saved vectors by applying $\Pi$, their
normal-word membership, and the separate raising premises remain
inputs from the authenticated producer records.  In particular,
this reader does not turn a stored projected vector into a fresh
calculation of its ambient vertex modes.

The $33$ projected vectors are also exported as a portable text
packet, \artifact{c8_reviewability_20260927/gate4_vectors.jsonl.gz}:
twelve targets followed by the $21$ normal columns.  Every line
contains its name, exact rational coefficients, and the digest
of its original vector record.  A header binds the separate
relation packet and records the original source hashes.
The packet contains $66610$ sparse state terms and is approximately
$2.45$ MB compressed.
The standalone verifier
\artifact{c8_reviewability_20260927/check_gate4_packet.py}
reads only this compressed text packet and
\artifact{c8_reviewability_20260927/finite_witnesses.json}.
It repeats all twelve full residual checks without opening an
original SQLite database.  Its output is
\artifact{c8_reviewability_20260927/gate4_packet_checked.json}.
All twelve residuals passed in $20.28$ seconds with peak RSS
$35380$ KiB.
This makes the residual arithmetic independently reproducible
from the companion bundle while preserving the preceding
boundary concerning the interpretation of the projected vectors.

The mathematical recovery step is short enough to verify directly.
The earlier exact raising controls give $E_Lt,E_Rt\in W_4$.
Since $W_4$ is stable under all current zero modes,
$[\Pi t]=[t]$ in $B_4/W_4$.  Since $N_4$ is current-stable,
the right side of \eqref{eq:reader-projected-residual} lies in
$N_4\subset W_4$.  Thus $t\in W_4$.  Neither this inference nor
the residual equality requires the $21$ chosen columns to span
the entire normal space.

To finish Gate 4, the saved matrices of $E_L,E_R$ on $P$ provide
$644$ columns, of rank $294$.  For each of the nine highest left
fields, append its certified seed columns; the augmented rank is
$322$.  The prime used to exhibit this full rank is
\[
 p=1000003,\qquad s\longmapsto231543,\qquad
                         t\longmapsto1000.
\]
Both defining quadratic equations hold modulo $p$, and all
denominators in a specialized column must be invertible.  A reader
can make the witness completely explicit by retaining the accepted
column indices and nonzero pivot factors during elimination: their
product, with the row-swap sign, is the determinant of the selected
minor.  The existing verifier recomputes the elimination; the new
packet does not claim to export or recheck these $322\times322$
minors.

The direction of the rank argument matters.  A nonzero minor after
specialization proves that its characteristic-zero determinant is
nonzero.  More precisely, the reduction is a ring homomorphism on
$\ZZ[D^{-1},s,t]\subset F$, where $D$ contains the finitely many
denominators used and $p\nmid D$; it is not a homomorphism from the
whole characteristic-zero field to $\mathbb F_p$.  Because $P$
already has exact dimension $322$, the nonzero minor proves
spanning of $P$.  It supplies no characteristic-zero upper bound
for an otherwise unknown space.  If $C_a$ is the kernel of
$[p]\mapsto[(a)_0p]$, the resulting equality is
$P=E_LP+E_RP+C_a$.  The kernel is raising-stable.  Cartan weights
on $P$ satisfy $|m_L|\leq3$, $|m_R|\leq21$; every raising step
increases $m_L+m_R$ by two.  Thus every raising word of length
$25$ vanishes.  Iterating the spanning equality in $P/C_a$ gives
$P/C_a=0$.  This is the finite step that recovers all $2898$
highest-left identities from the smaller collection of witnesses.

\subsection{Gate 5: the entire small scalar calculation}

Assume the already proved normal controls
$(S_2)_0N_4\subset N_5$ and $(N_3)_0S_3\subset N_5$.
If $P_4\subset N_4$, Gate 5 follows immediately.  Otherwise
$Q=W_4/N_4$ is the triplet with weights $(0,-2),(0,0),(0,2)$.
For its lowest vector $q_-$, define
\[
 \beta(a,p)=[a_0p]\in Q,\qquad
 F(a,q)=[a_0\widetilde q]\in B_5/N_5.
\]
The normal controls make $F$ well defined.  The two saved bottom
relations are
\begin{align*}
 \beta(a_3,p_3)&=c_3q_-,& p_3&=\texttt{r3\_000910},\\
 \beta(a_6,p_6)&=c_6q_-,& p_6&=\texttt{r3\_000899},\\
 c_3&=\frac{-593619+1349829s+1399567t-1072149st}{9831968},\\
 c_6&=-\frac{57967}{42018}+\frac{22295}{798342}st.
\end{align*}
Their input weights are $(0,-14)$ and $(-1,-9)$, respectively.
The vector $q_-$ is the image of the actual stored field
\texttt{P4\_minus}; this specifies its normalization.  Both
$c_3$ and $c_6$ are visibly nonzero, since $1,s,t,st$ are a
$\QQ$-basis and their rational coordinates are nonzero.
The equations are quotient identities: the full records also
give the normal-space terms that were suppressed here.

The lower-closure database reduces eight brackets.  Seven are
exactly zero as vectors in $S_3$:
\begin{equation}\label{eq:reader-seven-zeros}
 (a_0)_0a_3=(a_1)_0a_3=(a_2)_0a_3=(a_3)_0a_6
          =(a_5)_0a_3=(a_6)_0a_3=(a_7)_0a_3=0.
\end{equation}
There is one nonzero normal-product identity.  Define
\begin{align*}
 R_1&=\frac{1+4st}{25}(E_L)_{-1}a_1,\\
 R_2&=\frac{81-15st}{71}
           \left((E_L)_{-1}a_2-\frac{-7+3st}{20}R_1\right).
\end{align*}
Then the identity, with its denominators cleared, is
\begin{equation}\label{eq:reader-small-bracket}
 554600(a_4)_0a_6
       =(-216348s+83596t)R_1+(34265s+103895t)R_2.
\end{equation}
Both $R_1,R_2$ belong to $N_3=K_{{\rm cur},-1}S_2+TS_2$.
Their saved row names are \texttt{r3\_000141} and
\texttt{r3\_000150}.  Their definitions follow directly from
the two checkpoint recipes and the normalization rule
\eqref{eq:reader-row-recipe}; thus the displayed normal terms can
be understood without treating a row identifier as an unexplained
new field.

The precise lower records and weight exclusions are:
\begin{center}
\begin{tabular}{ccrcl}
\toprule
$i$&$j$&Lower record&$\operatorname{wt}((a_i)_0p_j)$&Bracket\\
\midrule
0&3&537&$(0,-10)$&zero\\
1&3&777&$(0,-6)$&zero\\
2&3&1017&$(0,-6)$&zero\\
3&6&1286&$(-1,3)$&zero\\
4&6&1526&$(0,-8)$&equation \eqref{eq:reader-small-bracket}\\
5&3&1737&$(1,-9)$&zero\\
6&3&1977&$(1,-7)$&zero\\
7&3&2217&$(1,-1)$&zero\\
\bottomrule
\end{tabular}
\end{center}
For direct comparison with those database rows, the primitive
right inputs have the one-term expansions
\[
 a_3=\frac{6+2s-4t+st}{14}\,\texttt{r2\_000047},
 \qquad
 a_6=\frac{307+297st}{1880}\,\texttt{r2\_000076}.
\]
For example, row $1526$ records the mode
$(a_4)_0\texttt{r2\_000076}$.  Multiplying its coordinates by
$(307+297st)/1880$ gives exactly
\eqref{eq:reader-small-bracket}.  The reader script independently
performs this rational recombination for all eight rows.  It
also checks their stored residuals and the two bottom records,
but does not recalculate these eight ambient mode products.

Every weight in the table is outside $Q$, so
$\beta(a_i,p_j)=0$.  The zero-mode Jacobi identity and the normal
controls give
\[
 F(a_i,\beta(a_j,p_j))=F(a_j,\beta(a_i,p_j))=0.
\]
The displayed nonzero scalars imply $F(a_i,q_-)=0$ for all eight
charged highest fields.  Raise this \emph{proved zero identity},
holding the left highest vector fixed, to obtain the remaining
right triplet inputs; then use current covariance for the left
descendants.  No assertion is made that the weight exclusion
continues to hold after raising $p_j$.

The retained neutral field requires one additional input: actual
charged brackets span $P$ as a current module.  Its certificate
uses $28$ charged seeds and $644$ lowering columns of rank $294$;
the augmented rank is $322$.  The same finite nilpotent argument
therefore gives the full current span.  Writing $e_0a_3$, modulo
$N_3$, as a sum of those charged brackets and applying Jacobi
twice proves the neutral bottom identity.  The other neutral
field follows by subtracting the affine stresses and $e$ from
$\omega_B$.  This completes the $24$ independent highest channels
as consequences of lower identities; it does not require $24$
separate new grade-five mode calculations.

\subsection{Exact location and scope of the witnesses}

The compact packet retains every coefficient in the twelve
projected identities and in the eight small lower combinations.
For comparison with the frozen sources, the following names are
unambiguous:
\begin{itemize}
\item The twelve full projected vectors and relation records are in
\artifact{hi_cft_20260923_c8two_h4_singlet_exact.sqlite}.
Their selected-column recipes are in
\artifact{hi_cft_20260923_c8two_h4_singlet_prime.sqlite}; the
underlying normal-word recipes are in
\artifact{hi_cft_20260923_next_n4_strong_skip_v2.sqlite}.
\item The separate raising premise is
\artifact{hi_cft_20260923_c8two_h4_singlet_premise_verified.json}.
All seed references, quotient coordinates and saturation inputs
are in
\artifact{hi_cft_20260923_c8two_structure_saturation.json}.
The completed mathematical reduction is recorded by
\artifact{hi_cft_20260923_c8two_gate4_verified.json}.
\item The eight rows in the table are primary keys of
\texttt{identities} in
\artifact{hi_cft_20260923_next_structure_lower_closure.sqlite}.
The two bottom relations are named
\texttt{A3\_zero\_r3\_000910} and
\texttt{A6\_zero\_r3\_000899} in
\artifact{hi_cft_20260923_c8two_targeted_h4_trimmed_v2/block_p0_m2.sqlite}.
\item The charged-span verifier output is
\artifact{hi_cft_20260923_c8two_charged_span_verified.json}.
The Gate 5 coefficient summary is
\artifact{hi_cft_20260923_c8two_gate5_verified.json}.
When replaying its original checker, the \texttt{--charged-span}
argument takes the producer
\artifact{hi_cft_20260923_c8two_structure_charged_span.json},
whose schema contains the actual matrices and seeds.
\end{itemize}

The independent reader's complete invocation is
\begin{verbatim}
python hi_cft_safe_run.py --limit-mib 1024 -- \
  python c8_reviewability_20260927/finite_witnesses.py \
  --output c8_reviewability_20260927/finite_witnesses_review.json
\end{verbatim}
For a bundle-only Gate 4 replay, the invocation is
\begin{verbatim}
python hi_cft_safe_run.py --limit-mib 1024 -- \
  python c8_reviewability_20260927/check_gate4_packet.py \
  --packet c8_reviewability_20260927/gate4_vectors.jsonl.gz \
  --relations c8_reviewability_20260927/finite_witnesses.json \
  --output /tmp/c8-gate4-packet-review.json
\end{verbatim}
Both checkers refuse to overwrite an output; choose another fresh name
if a review output already exists.  The SQLite reader requires its
output to remain inside the repository, whereas the portable checker
accepts the displayed temporary path.  They must be run by the coordinating
agent when the numerical slot is free.  Their allocation reviews bound
individual decoded records and retain only one target residual and
one source vector at a time.  The permanent launcher supplies the
desktop memory reserve, process-group watchdog, and two-CPU limit.

The source ledger and the small reader have complementary scopes.
The ledger authenticates the unchanged construction files and
fresh replay outputs.  The reader exposes selected mathematical
content and repeats a specified set of exact equalities.  The
finite-to-all-weights theorem still uses the proved filtered Wick
and recursive zero-product argument, together with the full lower
controls and Gates 4--5.  Printing a sample scalar relation would
not replace that theorem, just as hashing all its inputs would
not prove a residual identity.

% Supplement to the manuscript's embedded structural-data certificate.
% Requires amsmath, amssymb, booktabs and hyperref/url; no research macros.
% The accompanying TSV files are direct transcriptions, not new verifications.
\section{Structural inputs in a form suitable for direct checking}
\label{hr:structural-inputs}

The manuscript's embedded structural-data certificate has four parts:
the Potts factor, the level-thirteen parafermion factor and positive
stress identity, the lattice--Potts core and its six spectral witnesses,
and irreducibility on five finite top spaces. This section replaces
short success flags by the equations they assert. It also distinguishes
an explicit saved witness from a producer's assertion that such a
witness was found. The formulas below do not certify their own
membership in the algebra generated by the original fields.

\subsection{Coordinates and the meaning of a check}

Use the manuscript's simple-root basis $\alpha_0,\ldots,\alpha_7$,
Gram matrix $G$, lattice cocycle $\epsilon$, and transported positive
form. Write
\[
 s^2=13,\qquad t^2=-3,\qquad
 K=\mathbb Q(s,t),\qquad
 \overline{s}=s,\quad\overline{t}=-t.
\]
A stored coefficient $(q_0,q_s,q_t,q_{st})$ means
$q_0+q_ss+q_tt+q_{st}st$. The state labelled by
$(\beta;(i_1,\ldots,i_r))$ is
\[
 h_{i_1\bmod8}(-1-\lfloor i_1/8\rfloor)\cdots
 h_{i_r\bmod8}(-1-\lfloor i_r/8\rfloor)e^\beta.
\]
Thus index $i+8$ means $h_i(-2)$, not a second coordinate
direction. Ordered oscillator lists represent commuting Heisenberg
creation operators. The coefficient tables accompanying this section
retain every rational coefficient and contain no floating-point
approximations.

Write $D=L_B(-1)$, $:ab:=a_{-1}b$, and
$Y(a,z)=\sum_n a_nz^{-n-1}$. A polynomial identity below is an
equality of actual lattice vectors in this basis. A check consists of
expanding both sides, collecting equal basis states, and checking that
each of the four rational components of every coefficient vanishes.
No numerical tolerance enters this definition.

\subsection{The Potts fields: a compact formula and two worked checks}

Set
\[
 \alpha=(0,0,0,1,1,0,-1,-1),\qquad
 \beta=(1,1,0,-1,-1,0,0,0),\qquad \gamma=\alpha+\beta.
\]
Their Gram matrix is $\left(\begin{smallmatrix}4&-2\\-2&4\end{smallmatrix}\right)$.
The complete Potts stress vector is
\begin{equation}
 p=\frac{h_\alpha(-1)^2+h_\beta(-1)^2+
                 h_\alpha(-1)h_\beta(-1)}{15}\mathbf1
     +\sum_{\delta=\pm\alpha,\pm\beta,\pm\gamma}p_\delta e^\delta,
 \label{hr:potts-field}
\end{equation}
where
\[
\begin{array}{c|rr}
 \delta&p_\delta&p_{-\delta}\\\hline
 \alpha&-(s+3t)/100&(-s+3t)/10\\
 \beta&(7+3st)/100&(7-3st)/100\\
 \gamma&(-s+3t)/100&-(s+3t)/10
\end{array}
\]
Let $A=\alpha-\beta$, $C=\alpha+2\beta$, $E=-2\alpha-\beta$,
and $X_\delta=5p_\delta e^\delta$. Its weight-three generator is
\begin{equation}
 w_P=:h_Ah_Ch_E:
 -9\bigl(:h_A(X_\gamma+X_{-\gamma}):
       +:h_C(X_\alpha+X_{-\alpha}):
       +:h_E(X_\beta+X_{-\beta}):\bigr).
 \label{hr:potts-primary}
\end{equation}
The triple product is the indicated nested Heisenberg product; its
factors commute. These formulas contain the complete fields, not only
their projections to a selected subspace. The independent readable-input
checker has compared their complete 27- and 92-term expansions with
the saved fields and obtained exact equality.

Here is a short check of $p_3p=2\mathbf1/5$. The Heisenberg part of
$p$ is $2/5$ times the rank-two plane's conformal vector, and its
contribution to $p_3p$ is $4/25$. Each root has squared norm four,
so $\epsilon(\delta,-\delta)=1$ in the specified cocycle. The table
gives
\[
 p_\alpha p_{-\alpha}
   =\frac{s^2-9t^2}{1000}=\frac1{25},\qquad
 p_\beta p_{-\beta}
   =\frac{49-9(st)^2}{10000}=\frac1{25},\qquad
 p_\gamma p_{-\gamma}=\frac1{25}.
\]
Both orders of each root pair occur. The total is therefore
$4/25+6/25=2/5$, giving central charge $4/5$.

The transported PCT factor on $e^\delta$ is
$g(\delta)=\prod_i(5,8,9,8,5,1,2,2)_i^{\delta_i}$.
For $\alpha,\beta,\gamma$ these factors are $10,1,10$.
For example $10\overline{p_\alpha}=p_{-\alpha}$; the other
two pairs work in the same way. The quadratic Heisenberg part is
PCT fixed. This verifies $\theta p=p$ directly.

The remaining field and membership checks have a precise finite form:
\begin{align*}
 p_1p&=2p,&p_2p&=0,&p_0p&=Dp,\\
 p_1w_P&=3w_P,&p_nw_P&=0\quad(n\geq2),&
 L_B(1)w_P&=0,\\
 p&\in S_2,&w_P&\in K_{-1}S_2+DS_2+\operatorname{span}(S_{2,0}S_2).
\end{align*}
In the last expression $S_{2,0}S_2$ means all $a_0b$ with
$a,b\in S_2$; one keeps total Cartan weight zero when checking $w_P$.
The saved exact primary has 92 basis terms, and $p$ has 27.
The original exact membership calculation reached rank 41 after
starting with a rank-40 descendant span and adding one zero-weight
pair. Its old record did not save a membership expression. The new
export described in Section~\ref{hr:membership-words} supplies an
explicit expression with 33 nonzero coefficients. The standard
lattice Potts identification applies after this membership and the
stated field identities are verified; the rank 41 alone does not
establish membership.

\subsection{Parafermion normalization and the complete polynomial rows}

The exact parafermion stress $z$ and unnormalized primary $w$ are
given by 64 and 342 basis terms, respectively. They are transcribed
in \path{structural_fields.tsv}. Put $L=z$ and $a=\sigma w$, where
\[
 \sigma^2=-\frac{7+3st}{9223500},\qquad
 w_5w=(-86086+36894st)\mathbf1.
\]
No choice of square root is needed in the stored even identities;
odd identities have a common factor $\sigma$. Indeed, $(st)^2=-39$
gives the worked normalization check
\[
 (-86086+36894st)(-7-3st)=4919200,
 \qquad \sigma^2w_5w=\frac8{15}\mathbf1.
\]
Define
\[
 c=\frac85,\qquad \ell=\frac{14}{473},\qquad
 b=a_1a,\quad d=a_1b,\quad R=a_1d.
\]
All normal products in the following formulas retain their indicated
order and parentheses. The full singular self-row is
\begin{equation}
\begin{array}{c|rrrrrr}
 n&5&4&3&2&1&0\\\hline
 a_na&\frac8{15}\mathbf1&0&2L&DL&b&
       \frac12Db-\frac1{12}D^3L .
\end{array}
\label{hr:pf-selfrow}
\end{equation}
Products with larger mode index vanish. The second row is
\begin{align}
 a_nb&=0\quad(n\geq4),\notag\\
 a_3b&=\frac{69283}{2365}a,\notag\\
 a_2b&=\frac83\bigl(5-2\ell(2+c)\bigr)Da,\notag\\
 a_1b&=d,\notag\\
 a_0b&=\frac25Dd+\frac{32\ell}{5}:LDa:
          -\frac{48\ell}{5}:DLa:
          +\frac2{15}\bigl(-5+2\ell(c-1)\bigr)D^3a.
 \label{hr:pf-ab}
\end{align}
For example the single coupling check is exactly
\[
 \frac{31-69283/2365}{16(2+8/5)}
 =\frac{4032/2365}{288/5}=\frac{14}{473}.
\]
This arithmetic is necessary but is not the parafermion
identification: the nonlinear rows and decoupling relation are needed.

The upper part of the third row is
\begin{align*}
 a_7d&=c\bigl(55-16\ell(2+c)\bigr)\mathbf1,\\
 a_6d&=0,\\
 a_5d&=\frac43\bigl(175-64\ell(2+c)\bigr)L,\\
 a_4d&=\bigl(110-32\ell(2+c)\bigr)DL,\\
 a_3d&=\bigl(95-16\ell(11+4c)\bigr)b
          +128\ell:LL:+\bigl(10+8\ell(c-1)\bigr)D^2L,\\
 a_2d&=64\ell:DLL:
          +\bigl(75/2-4\ell(13+5c)\bigr)Db
          +\frac{-25+8\ell(9+5c)}{12}D^3L.
\end{align*}
Here $:DLL:$ means $:(DL)L:$, not $D(:LL:)$.
For $a_1d=R$ the exact polynomial is
\begin{align}
 R={}&-\frac{148744}{473}:Lb:
 +\frac{566160}{473}:aa:
 -\frac{4800}{121}:L(:LL:):\notag\\
 &+\frac{3860}{11}:(D^2L)L:
 -\frac{33320}{121}:(DL)(DL):
 -\frac{489817}{4730}D^2b
 +\frac{2193061}{52030}D^4L.
 \label{hr:pf-decoupling}
\end{align}
The last simple pole is
\begin{align}
 a_0d={}&\frac13DR+\frac{448}{1419}:LDb:
 -\frac{896}{1419}:(DL)b:
 -\frac{224}{1419}:(D^3L)L:\notag\\
 &-\frac{34019}{28380}D^3b
 +\frac{9473}{170280}D^5L.
 \label{hr:pf-last-pole}
\end{align}
Thus every entry of all three singular rows is an explicit
differential normal polynomial in $L,a,b,d$.

There is a useful small algebraic alternative to directly expanding
the weight-seven vector in \eqref{hr:pf-last-pole}. Write
$a_0b=\beta Dd+Q$, where $\beta=2/5$ and $Q$ is the last three
terms of \eqref{hr:pf-ab}. The iterate identity and self-skew symmetry give
\begin{equation}
 (1-\beta)a_0d=\beta DR+a_1Q
       -\frac14D(b_1b)+\frac1{48}D^3(b_3b)
       -\frac1{480}D^5(b_5b).
 \label{hr:pf-derived-pole}
\end{equation}
The terms $b_1b,b_3b,b_5b$ are obtained from the already given
upper poles and $a_0b$, without using $a_0d$. Reordering uses, for
example,
\[
 :aDa:-:(Da)a:=\frac1{12}D^3b-\frac1{60}D^5L,
 \qquad :L(D^2L):-:(D^2L)L:=\frac16D^4L.
\]
Substituting $c,\ell$ in \eqref{hr:pf-derived-pole} gives precisely
\eqref{hr:pf-last-pole}. A scalar verification of this implication
checks an algebraic consequence of the other field identities; it
does not independently verify those identities in the lattice VOA.

The seven coefficients of \eqref{hr:pf-decoupling} also agree with
the Dong--Lam--Yamada parafermion formulas after the recorded change
of generators. The manuscript's finite reconstruction proposition
then transports the full polynomial presentation to the actual
simple generated algebra. The required actual-field inputs include
\[
 z_1z=2z,\quad z_2z=0,\quad z_3z=\tfrac45\mathbf1,
 \quad z_1w=3w,\quad z_nw=0\ (n\geq2),\quad z_0w=Dw,
\]
PCT stability, and $z\in S_2$, $w\in S_3$. The membership
construction reaches the 58-dimensional Cartan-zero grade-three
span after 154 zero-weight products. Its new export supplies an
expression for $w$ with 42 nonzero coefficients, as described next.
The recorded computation of the full nonlinear product $a_0b$
involves 25,045 nonzero basis terms; $R$ has 28,342.
These are the substantial vector identities behind the short
polynomials above. They can be checked by a sparse exact replay,
but are not replaced by the two worked scalar checks.

\subsection{Explicit membership expressions in generated words}
\label{hr:membership-words}

The new file \path{membership_witnesses.json} contains 120
explicit generated $S_2$ vectors $v_j$, 58 independent raw
grade-three source vectors $g_i$, and the two identities
\[
 w_P=\sum_{i\in I_P}\mu_i g_i,\quad |I_P|=33,
 \qquad
 w=\sum_{i\in I_{13}}\nu_i g_i,\quad |I_{13}|=42.
\]
Every $g_i$ is labelled as $J_{-1}v_j$, $Dv_j$, or
$(v_j)_0v_k$. The complete coefficient lists are in
\path{membership_witnesses.coefficients.tsv}; the JSON also
contains every source and target coordinate. Thus the proposed
linear identities are independently checkable without repeating
the search for a suitable span or choosing new pivots.

For example, in the Potts expression
\[
 \mu_0=\frac{1053}{70},\qquad
 \mu_1=-\frac{729}{280},\qquad
 \mu_2=\frac{1647}{280},\qquad
 \mu_{40}=-\frac1{3985293312000},
\]
where
\begin{align*}
 g_0&=(H_\mathrm{left})_{-1}v_0,
 &v_0&=(H_\mathrm{left})_{-2}\mathbf1,\\
 g_1&=(H_\mathrm{left})_{-1}v_1,
 &v_1&=(H_\mathrm{left})_{-1}H_\mathrm{left},\\
 g_2&=(H_\mathrm{left})_{-1}v_4,
 &v_4&=(H_\mathrm{left})_{-1}H_\mathrm{right},\\
 g_{40}&=(v_{92})_0v_{119}.&&
\end{align*}
The remaining 29 coefficients and the exact $v_{92},v_{119}$
appear in the same files. These are raw source words, not normalized
row-reduction vectors with an unstated change of basis.

The exporter uses the original low-grade mode engine and tracks all
linear combinations. It finished with 43,462 persistent sparse terms
and about 83 MiB peak resident memory. An independent coordinate-sum
checker, \path{check_membership_coordinates.py}, uses only exact
rational arithmetic in $K$. Its passing record checks both literal
sums and all 21,772 decoded basis rows, including source grades and
the Cartan-zero condition; duplicate basis rows and duplicate JSON
keys are rejected. Two boundaries remain separate:
checking a sum of supplied vectors is not checking that each
supplied vector equals the mode word named by its recipe, and the
generation of all 120 $v_j$ by the original short generating set
uses the finite construction certificate. The export makes both
boundaries explicit instead of encoding membership by a truth value.

\subsection{Where the positive stress identity is checked}

Section~\ref{hr:positive-stress-check} prints the full four-axis
Gram and overlap table and checks the resulting rational identity.
The independent calculation uses the actual saved $p,z,e$ and the
affine stresses; it imports no original mode engine. Its passing
record is \path{c8_reviewability_20260927/readable_checks.json},
and the exact fields are in \path{grade2_fields.tsv}.
Identifying these orthogonal projections with Haar averages still
uses the certified fact that the four stresses span the
current-invariant part of the constructed $S_2$. This is a statement
about $S_2$, and does not require $S_2=U_2$ or properness.

\subsection{The lattice core and six small spectral witnesses}

Return to the lattice vectors $\ell,r$ from the small lattice
calculation, and write $H=r$ for the right Cartan vector.  The rank-two
lattice core has basis $\ell$ and $v=(1,1,0,1,1,-3,1,1)$, with
\[
 G_{\mathrm{core}}=\begin{pmatrix}2&3\\3&24\end{pmatrix},
 \quad\det G_{\mathrm{core}}=48-9=39,
 \quad \rho=v/3=\ell/2+H/6,
 \quad \rho^2/2=4/3.
\]
The actual three-term field $u$ supplies the compact identities
\[
 p_1u=\tfrac23u,\qquad p_2u=p_3u=0,
 \qquad u_{-5}(u_{-3}u)=-\tfrac35e^v.
\]
The last is a leading-coefficient lattice identity. Its small
coefficient can be checked directly from the three exponential
terms of $u$; it does not require constructing an ambient basis in
weight twelve. Together with PCT and the left currents, these
identities place the rank-two lattice and the charged simple-current
sectors in the generated algebra.

The resulting core stress is
\[
 \omega_{\mathcal T}
 =\tfrac14h_{\ell}(-1)^2+\tfrac1{156}h_H(-1)^2+p,
 \qquad c_{\mathcal T}=14/5.
\]
Put $q=\omega_B-\omega_{\mathcal T}-z$, of central charge $18/5$.
The six order-thirteen witnesses are as follows. Every row names a
nonzero actual vector in $S_2$; each corresponding Cartan block
inside the constructed $S_2$ has dimension three.
\begin{equation}
\begin{array}{c|rrrrr}
 (h_{\ell}(0),h_H(0))&\dim\ker(q_1)&q_1&z_1&p_1&
                      (\omega_{\mathcal T})_1\\\hline
 (-1,-5)&1&0&12/13&2/3&14/13\\
 (-1, 5)&1&0&12/13&2/3&14/13\\
 ( 0,-8)&1&0&12/13&2/3&14/13\\
 ( 0, 8)&1&0&12/13&2/3&14/13\\
 ( 1,-5)&1&0&12/13&2/3&14/13\\
 ( 1, 5)&1&0&12/13&2/3&14/13
\end{array}
\label{hr:core-eigenvalues}
\end{equation}
The full six vectors, each with 21 basis terms, are printed in
\path{core_spectral_witnesses.tsv}. For example, the first vector
has nonzero coefficient
$35/16+(175/1872)st$ at momentum $(-1,0,1,1,1,1,0,0)$,
and coefficient $-203s/208-161t/48$ at $(0,0,0,0,0,1,1,1)$.
The former cannot vanish, since $(st)^2=-39$ while its putative
rational value $-234/10$ would have positive square.
Each line of \eqref{hr:core-eigenvalues} asserts the residual equations
\[
 q_1v=0,
 \quad(13z_1-12)v=0,
 \quad(3p_1-2)v=0,
 \quad(13(\omega_{\mathcal T})_1-14)v=0,
 \quad v\in S_2.
\]
These can be checked separately, one 21-term vector at a time;
the table of eigenvalues is not by itself the verification.

For a hand check of the module label, parafermion lowest weights are
\[
 h_{i,j}=\frac{i(13-i)}{390}+\frac{j(i-j)}{13},
 \qquad 1\leq i\leq13,\quad 0\leq j<i.
\]
Solving $h_{i,j}=12/13$ means
$i(13-i)+30j(i-j)=360$. Reduction modulo three forces
$i\equiv0,1\pmod3$. The candidates $i=1,3,4,6,7,9,10,12$
give, respectively, $i(13-i)=12,30,36,42,42,36,30,12$.
Except for $i=3,10$, the difference from 360 is not divisible
by 30. For $i=3,10$ it would require $j(i-j)=11$, impossible
for the permitted integral $j$. Thus $i=13$ and
$j(13-j)=12$, giving $j=1,12$. These are inverse simple
currents. This finite arithmetic check explains why the specific
energy $12/13$ matters; it does not produce the lattice eigenvectors.

\subsection{Finite irreducibility: the criterion and actual row recipes}

The following elementary argument turns cyclicity into the precise
irreducibility needed here. Let $E$ be a finite-dimensional Hilbert
space and $\mathcal A\subset\operatorname{End}(E)$ an adjoint-closed
algebra containing two commuting selfadjoint Cartan operators.
Suppose a joint eigenspace is the line $\mathbb Cv$, and
$\mathcal Av=E$. If $W$ is invariant, its orthogonal complement
is invariant, so the orthogonal projection $P_W$ commutes with
$\mathcal A$. It preserves $\mathbb Cv$ and acts there as zero
or one. Commuting with every $A\in\mathcal A$ and using cyclicity
then gives $P_W=0$ or $1$. Hence $E$ is irreducible.

Let $K_{\rm cur}$ be the span of the six affine currents.
For the eighty-current space, $E=K_{\rm cur}^\perp\cap B_1$ has
dimension 80. The operators are current zero modes and the zero modes $a_1$
of the named weight-two fields and their adjoints. The distinguished
Cartan line has weight $(0,10)$. The old v2 record saved only
assertions and dimensions. The new
\path{current80_witness.json} instead gives the full 86-state
ambient basis and positive Hermitian Gram matrix, an explicit
80-row basis of $E$, all 30 zero-mode and adjoint matrices,
and all 80 cyclic rows with their parent, operator, elimination
coefficients, normalization and pivot. The matrices have only
3,486 nonzero entries altogether. The exporter recomputes
invariance, adjointness and lowering-mode annihilation with the
original exact low-grade engine; it took about 1.3 seconds and
36 MiB. Thus these are finite matrices and row identities available
for inspection, rather than an unexplained rank assertion.

A first cyclicity step can also be checked directly. The seed is
$e^\beta$, where $\beta=(1,1,1,1,1,0,0,0)$. The left and right
Cartan weights are $0$ and $10$. In the prescribed cocycle,
\[
 (F_\mathrm{right})_0e^\beta
 =5e^{(0,1,1,1,1,0,0,0)}
  -5e^{(1,1,1,1,0,0,0,0)}.
\]
Only the negative simple roots $-\alpha_0,-\alpha_4$ contribute:
their pairings with $\beta$ are $-1$, and the corresponding
cocycle signs are $-1,+1$. Both coefficients in $F_\mathrm{right}$
are $-5$. Dividing this displayed result by five gives the first
nontrivial normalized cyclic row in the export.

The independent stdlib checker
\path{check_current80_matrices.py} has passed; its result is
\path{current80_matrix_checks.json}. It checks the lattice Gram matrix
by exact positive $LDL^{\mathsf T}$ decomposition, the root-state
metric by the positive torus factors, all fifteen adjoint pairs,
orthogonality and independence of the complement, invariance under
all thirty matrices, the two Cartan matrices and their unique
$(0,10)$ eigenline, and the eighty cyclic row recipes. The eight
positive diagonal entries of $D$ are
\[
 2,\quad\frac32,\quad\frac43,\quad\frac54,\quad
 \frac65,\quad\frac12,\quad2,\quad\frac32.
\]
The check used about 25 MiB of resident memory. Its scope is
the supplied finite matrices. Identifying those matrices with the
actual field modes and checking lowering-mode annihilation remain
separate from this independent linear-algebra verification.

The four 81-dimensional mixed tops have more complete saved data.
Each has 81 root states of squared norm two, a positive diagonal
metric, all 81 elimination rows, and their construction recipes.
All four have the same Cartan multiplicities:
\[
\begin{array}{c|rrrrrrrr}
 h_{\ell}(0)=\pm1,\ h_H(0)&-7&-5&-3&-1&1&3&5&7\\\hline
 \text{multiplicity}&1&2&3&3&3&3&2&1
\end{array}
\]
\[
\begin{array}{c|rrrrrrrrrrr}
 h_{\ell}(0)=0,\ h_H(0)&-10&-8&-6&-4&-2&0&2&4&6&8&10\\\hline
 \text{multiplicity}&1&2&4&5&7&7&7&5&4&2&1 .
\end{array}
\]
Each outer row sums to 18 and the middle row to 45, giving
$18+45+18=81$. In particular $(0,10)$ is a unique line.
The saved seed has index 68 in the ordered root list. Multiplying
its momentum by three gives
\[
\begin{array}{c|c}
 (\text{glue sign},\text{sector sign})&3\beta_{68}\\\hline
 (+,+)&(1,2,3,4,2,0,2,1)\\
 (+,-)&(2,4,3,2,1,0,1,2)\\
 (-,+)&(1,2,3,4,2,0,1,2)\\
 (-,-)&(2,4,3,2,1,0,2,1)
\end{array}
\]

Here is the complete first nontrivial cyclicity step in the $(+,+)$
case. Denote the root basis by $e_j$ and let $M_9$ be the adjoint
of the right raising-current zero mode, in the exact metric.
The saved recipe is
\[
 r_0=e_{68},\qquad
 r_1=\tfrac18M_9r_0=e_{62}-\tfrac14e_{67}.
\]
Thus the actual matrix identity to check is
$M_9e_{68}=8e_{62}-2e_{67}$, with
\[
 3\beta_{62}=(1,2,3,1,2,0,2,1),\qquad
 3\beta_{67}=(1,2,3,4,2,0,-1,1).
\]
Both output states have Cartan weight $(0,8)$, as predicted by
the current's weight. Subsequent steps have the exact format
\begin{equation}
 r_k=\lambda_k\left(M_{j_k}r_{p_k}
                   -\sum_{i<k}c_{ki}r_i\right),
 \qquad \lambda_k\ne0,
 \label{hr:cyclic-recipes}
\end{equation}
and each $r_k$ has a new normalized pivot. The saved recipes list
$p_k,j_k,\lambda_k,c_{ki}$ explicitly; all $p_k<k$. Checking
\eqref{hr:cyclic-recipes} and distinct pivots proves that 81
independent vectors belong to the cyclic span. Counting the saved
rows without checking their recipes does not establish cyclicity.

The matrices themselves need not be reconstructed by a general VOA
engine. On a pure root state $e^\beta$, the relevant zero modes of
the weight-one and weight-two summands have the elementary rules
\[
 \begin{array}{c|c}
 \text{state}&\text{zero-mode coefficient or action}\\\hline
 h_i(-1)\mathbf1&(\alpha_i,\beta)\\
 h_i(-2)\mathbf1&-(\alpha_i,\beta)\\
 h_i(-1)h_j(-1)\mathbf1&(\alpha_i,\beta)(\alpha_j,\beta)\\
 e^\alpha&g_0(\alpha)\epsilon'(\alpha,\beta)e^{\alpha+\beta}\\
 h_i(-1)e^\alpha&(\alpha_i,\beta)g_0(\alpha)
                          \epsilon'(\alpha,\beta)e^{\alpha+\beta}.
 \end{array}
\]
In the last two rows the term contributes only when
$\alpha+\beta$ is in the same norm-two root list.
The cocycle $\epsilon'$ of the chosen overlattice and its quadratic
gauge correction $g_0$ must be used; they are distinct from the
positive torus factor in the Hermitian metric. With diagonal metric
$\operatorname{diag}(d_i)$, each adjoint entry is checked by
\[
 (M^*)_{ji}=\overline{M_{ij}}\frac{d_i}{d_j}.
\]
This is a finite exact matrix problem. Passing from these irreducible
tops to entire irreducible restrictions uses the independently
proved bound $h<2$ and regularity. It is not a consequence of
finite cyclicity alone, and the four mixed restrictions are not
asserted to be mutually isomorphic.

\subsection{Source-to-equation map and remaining verification boundary}

Paths beginning \texttt{hi\_cft\_} refer to the repository root.
Unprefixed companion filenames throughout this section, such as
\path{structural_fields.tsv} and \path{current80_witness.json}, are
relative to \path{c8_reviewability_20260927/}; explicitly prefixed
paths remain repository-relative.  The accompanying TSV files
contain exact transcriptions, not new checks.
\begin{itemize}
\item \path{hi_cft_20260923_proof_potts_exact.json}:
keys \texttt{potts\_stress}, \texttt{weight\_three\_primary};
the compact formulas \eqref{hr:potts-field}--\eqref{hr:potts-primary}.
The earlier \path{hi_cft_20260923_proof_potts.json} contains the
exact stress checks but only a numerical test of primary membership;
the later exact record supersedes that numerical membership test.
\item \path{hi_cft_20260923_proof_pf13_exact.json}:
keys \texttt{stress}, \texttt{primary} and
\path{K_invariant_stress_projection_coefficients}.
Membership, PCT and translation use the separate
\path{hi_cft_20260923_proof_pf13_membership.json} and
\path{hi_cft_20260923_proof_remaining_coset.json}.
\item The parafermion rows use
\path{hi_cft_20260923_proof_pf13_upper_poles.json},
\path{hi_cft_20260923_proof_pf13_nonlinear.json},
\path{hi_cft_20260923_proof_pf13_truncation.json},
\path{hi_cft_20260923_proof_pf13_simple_pole_algebra.json}, and
\path{hi_cft_20260923_proof_pf13_dly.json}, with the producer
formulas specified by those records. The upper-pole continuation
imports fifteen completed assertions from the earlier producer log;
its final record is not sixteen freshly independent calculations.
\item \path{hi_cft_20260923_next_structure_core_extension.json}:
for each entry of \texttt{blocks}, use its
\texttt{Cartan\_weight}, then the \texttt{joint\_eigenspaces}
entry with \texttt{PF13\_weight=12/13}, then
\texttt{vectors[0].vector}. These are the six full witnesses
in \path{core_spectral_witnesses.tsv}.
\item \path{hi_cft_20260923_proof_current_module_v2.json}:
the original eighty-current assertions. These are supplemented by
\path{current80_witness.json}, containing explicit matrices and
cyclic-basis recipes, and the separate independent matrix checker.
\item \path{hi_cft_20260923_next_81top_all.json}:
each entry of \texttt{sectors} contains \texttt{roots},
\texttt{metric}, \texttt{rows}, and \texttt{recipes}.
Its \texttt{operator\_matrices} entries contain hashes and sparse
term counts, not the matrix entries themselves. The elementary
formulas above reconstruct the latter. The files
\path{top81_seeds_and_first_steps.tsv} and
\path{top81_pivots.tsv} make the small identifying data readable.
\end{itemize}

The most economical further certificates would therefore be:
independent checks of the exported membership words' mode-product
provenance and the finite top matrices' actual field actions, and independent
coefficient-by-coefficient checks of the substantial parafermion
normal-product identities. The small scalar and matrix checks printed
here expose their mathematical content and allow parts to be
verified by hand. They do not remove those larger exact-vector
obligations or formalize the ambient VOA in a proof assistant.


\section{Reproduction and the remaining large checks}
The new small checker is
\artifact{c8_reviewability_20260927/check_readable_inputs.py}. It uses
only Python's exact rational arithmetic and the displayed formulas;
it imports none of the research mode engine. Run from the repository:
\begin{quote}\small\begin{verbatim}
python hi_cft_safe_run.py --limit-mib 512 -- \
  timeout 120s python \
  c8_reviewability_20260927/check_readable_inputs.py \
  --output-dir /tmp/c8-readable-review
\end{verbatim}\end{quote}
This writes a fresh report in the chosen output directory; use a new
directory if it already exists. The report from this revision is
\artifact{c8_reviewability_20260927/readable_checks.json}.
The source hashes in that report identify the exact coordinate inputs.
The mathematical content is the recomputation; the hashes identify
which inputs were used. The complete current singular brackets,
the highest-vector conditions, the stress Gram matrix, the overlaps,
the projected identity, both complete compact Potts formulas and the
eight primary energy sums are checked.
The script does not check every higher OPE or prove that a saved field
belongs to the algebra generated by the short defining set.

The old large calculations remain relevant. In particular, dimensions
$120$, $956$ and $322$ and the full Gate~4 closure are not established
by inspecting a representative row. The preceding sections specify
which actual saved witnesses can be checked and which missing witness
formats would make the remaining verification easier. The companion
\artifact{hi_cft_c8_lean_assessment.pdf} explains how these finite
interfaces could become kernel-checked theorems, and why doing so would
still leave substantial VOA, tensor-category and analytic foundations.
\end{document}
